\chapter{Идея целиком}

\section{Три плоскости}

Проще всего увидеть архитектуру как три плоскости, которые в обычной агентной системе склеены в одном процессе, а здесь разведены по разным машинам и разным владельцам.

\begin{center}
\begin{tikzpicture}[node distance=6mm]
\node[boxa,minimum width=134mm,minimum height=13mm] (think) {Плоскость мышления\\[1pt]\footnotesize агентские циклы на машинах контрибьюторов: модель, контекст, ключи, ёмкость};
\node[boxg,minimum width=134mm,minimum height=13mm,below=of think] (ctrl) {Плоскость управления\\[1pt]\footnotesize кор-ноды: журнал записей, правила, полномочия, оркестратор};
\node[boxw,minimum width=134mm,minimum height=13mm,below=of ctrl] (exec) {Плоскость исполнения\\[1pt]\footnotesize среда проекта: рабочее дерево, санбоксы, тесты, основная ветка};
\draw[darr] (think) -- node[lbl,right] {назначения, ревизии, вердикты} (ctrl);
\draw[arr] (ctrl) -- node[lbl,right] {записи с сертификатом кворума} (exec);
\draw[arr] ($(think.west)+(-4mm,0)$) -- ++(-6mm,0) |- node[lbl,left,pos=0.25,rotate=90,anchor=south] {удалённый shell в санбоксе задачи} ($(exec.west)+(-4mm,0)$);
\end{tikzpicture}
\end{center}

Плоскость мышления это агентские циклы: модель, её контекст, инструменты, которые она вызывает. Они работают на машинах контрибьюторов под их собственными учётными данными и с их собственной ёмкостью. Плоскость исполнения это среда проекта: единственная машина или небольшой пул машин, где лежит код, собираются артефакты и гоняются тесты. Плоскость управления это журнал записей и кор-ноды, которые его ведут: кто что вправе делать, что назначено, что проверено, что интегрировано.

Между плоскостями проходят только команды, их вывод, хеши и подписи. Учётные данные не покидают машину контрибьютора. Код не покидает среду проекта иначе как в санбоксе конкретной задачи. Полномочия существуют только в журнале.

\section{Инверсия: среда одна, агенты снаружи}

Обычная схема: агент запускается там, где код. Наша: код лежит в одном месте, а агенты подключаются к нему снаружи. Среда не запускает агентских циклов вообще. Она пассивна: исполняет команды, поступающие извне, и только в границах, заданных журналом.

Это инверсия направления, и у неё есть цена и выигрыш. Цена: среда становится единственной точкой, через которую проходит вся работа, и её ёмкость ограничивает размер роя. Выигрыш: среду не нужно воспроизводить, ключи не нужно раздавать, а квоты приносят сами участники. Для дорогих сред выигрыш перевешивает.

\begin{whybox}[Почему не облачная песочница на агента]
Облачные песочницы решают изоляцию, воспроизводя среду на каждого агента. Для Xcode, ключей подписи, устройств, GPU и данных с ограничением на вывоз это либо невозможно, либо дорого. Кроме того, песочница всё равно принадлежит одному аккаунту с одной квотой; объединение квот разных людей она не решает.
\end{whybox}

\section{Журнал полномочий вместо git как шины}

Git остаётся в системе на двух ролях: хранилище, адресуемое по содержимому (любой коммит имеет хеш, и по хешу можно проверить, что именно проверял ревьювер), и механизм ветвления внутри одной машины. Git перестаёт быть шиной координации между участниками. Координация идёт через журнал.

Журнал это упорядоченная цепочка подписанных записей: цель поставлена, задача создана, задача назначена такой-то ноде, ревизия с таким-то хешем представлена, ревью такое-то, тесты такие-то, интеграция такая-то. Каждая запись ссылается на предыдущую по хешу и подписана достаточным числом кор-нод. Среда исполняет только то, что подтверждено журналом: открывает санбокс по записи назначения, сливает в основную ветку по записи об интеграции.

Так всякое действие в системе имеет ответ на вопрос кто и на каком основании его разрешил, и ответ этот проверяем без доверия к тому, кто его даёт.

\section{Сессия это кэш, журнал это истина}

У агента есть контекст: история диалога, промежуточные выводы, план. В обычных системах этот контекст и есть состояние проекта: если процесс упал, состояние потеряно. Здесь состояние проекта живёт в журнале, а контекст модели есть кэш. Оркестратор, который потерял сессию, восстанавливается из журнала. Оркестратор, который получил сессию от предшественника, обязан считать её недоверенной и разрешать любое расхождение в пользу журнала.

Это одно из самых важных решений в проекте, и глава 7 объясняет, почему без него не работает ни смена оркестратора, ни защита от отравленного контекста.

\section{Что мы не верифицируем}

Через консенсус кор-нод проходят только детерминированные предикаты: допустим ли переход состояния, лежит ли область задачи внутри области цели, относятся ли ревью и тесты к тому же хешу, что интегрируется, выбран ли ревьювер по случайности, которую никто не контролировал. Кор-ноды не голосуют о том, хорош ли код и прав ли ревьювер.

Это принципиально. Криптография и консенсус умеют доказывать, что процедура соблюдена: кто, когда и на каком основании. Они не умеют доказывать, что модель правильно поняла задачу. Попытка проголосовать о правильности суждения ломает и консенсус, и суждение. Поэтому суждение остаётся снаружи, а протокол делает его независимым (продьюсер, ревьювер и тестировщик разные), случайным (ревьювера нельзя подобрать) и подотчётным (каждое действие оставляет след).

\section{Уровни доверия}

Протокол определён для трёх уровней доверия с одинаковыми инвариантами и разным определением того, сколько подписей нужно, чтобы запись считалась совершённой.

\begin{center}
\small
\begin{tabular}{p{1.5cm} p{3.3cm} p{3.1cm} p{5.2cm}}
\toprule
Уровень & Сколько подписей & Порядок записей & От чего защищает сертификат \\
\midrule
L0 & одна, единственная кор-нода & единственная нода & от действий вне кора; кор и среда доверены \\
L1 & больше половины из $n$ & доверенный секвенсер & от одиночного несанкционированного решения при компрометации меньшинства ключей; от раздвоения истории защищает только честность секвенсера \\
L2 & $2f+1$ из $n \geq 3f+1$ & византийски устойчивый консенсус & и от несанкционированного решения, и от раздвоения истории при $f$ злонамеренных кор-нодах \\
\bottomrule
\end{tabular}
\end{center}

Первое развёртывание идёт на L0: одна кор-нода, один подписант. Все механизмы при этом уже работают, просто кворум равен одной подписи. Повышение уровня не требует переписывать ничего, кроме определения кворума. Последний столбец важнее, чем кажется: читатель должен понимать, защищает ли сертификат от одиночного несанкционированного решения, от раздвоения истории или от обоих. Глава 3 объясняет, откуда берутся $2f+1$ и $3f+1$ и почему порог большинства на L1 не даёт единственности истории.

\section{Десять инвариантов}

Всё остальное в книге есть способы обеспечить десять инвариантов. Здесь они перечислены кратко, чтобы читатель видел каркас; глава 5 разбирает каждый и показывает, что ломается без него.

\begin{enumerate}[label=I\arabic*.,leftmargin=2.6em]
\item Ни одного эффекта на среде без совершённой и актуальной записи авторизации с совпадающим digest плана.
\item Полномочия только сужаются: от цели к задаче и от задачи к ревизии; условия приёмки меняются только новой версией контракта с авторизацией.
\item Основная ветка продвигается только атомарной заменой на неизменяемого кандидата, к которому адресованы финальное одобрение и аттестации.
\item Суждение не входит в консенсус.
\item Продьюсер, ревьювер и исполнитель проверок независимы согласно объявленной политике, которая не ослабляется автоматически; случайность выбора фиксируется после того, как подбирать уже нечего.
\item Ключи контрибьютора не покидают его машину; секреты проекта существуют только в раннере.
\item Журнал есть истина, сессия есть кэш; записи устаревшей аренды отвергаются.
\item Любая кор-нода останавливает, кворум возобновляет; откат журнала невозможен.
\item Ограничение есть сообщаемый факт; требующий полного ограничения отказывает на частичном.
\item Все обращения к среде через протокол имеют проверяемое происхождение; полнота внутреннего контекста недоверенной ноды не удостоверяется.
\end{enumerate}

Если вы запомните из книги только этот список и три плоскости, вы уже сможете отличить правильное решение от неправильного почти в любом вопросе реализации.
